\section{Experimental Setup}
\label{sec:setup}

\subsection{Models}

The programme spans five model families and eight distinct checkpoints, all frozen; no weights are updated anywhere. The steering doctrine is developed and confirmed on three plants from two families: Qwen3-4B is the primary plant for cross-plant transfer, the lens benchmark, and comparative evaluation \cite{Qwen3_2025}; Qwen3.6-27B, from the same lineage at 6.75 times the parameter count, is used for the loadability scaling comparison; and Nemotron-Mini-4B, a pruned-and-distilled compact model \cite{Minitron2024} that also serves as the predictive-world-model twin in the author's related stack, is the plant on which the doctrine was first developed. The recognition-gated hardening study (Section~\ref{sec:results:moat}) extends coverage to four further compact instruct models across three additional families: Gemma-2-2B \cite{Gemma2_2024}, Llama-3.2-1B \cite{Llama3_2024}, Qwen2.5-1.5B \cite{Qwen25_2024}, and SmolLM2-1.7B \cite{SmolLM2_2025}. Table~\ref{tab:models} lists the full set with the loops in which each is used.

\begin{table}[t]
\caption{Model families and checkpoints across the programme. Layer counts and band ranges are the characterised values where measured.}
\label{tab:models}
\vskip 0.1in
\centering
{\footnotesize
\begin{tabular}{llll}
\toprule
\textbf{Family} & \textbf{Checkpoint} & \textbf{Size} & \textbf{Role} \\
\midrule
Qwen3     & Qwen3-4B          & 4B   & primary plant, L1--L22 \\
Qwen3     & Qwen3.6-27B       & 27B  & loadability scaling, L1b \\
Nemotron  & Nemotron-Mini-4B  & 4B   & doctrine development, twin \\
Gemma-2   & Gemma-2-2B        & 2B   & moat, works (L26) \\
Llama-3.2 & Llama-3.2-1B      & 1B   & moat, works (L26) \\
Qwen2.5   & Qwen2.5-1.5B      & 1.5B & moat, fails (L26) \\
SmolLM2   & SmolLM2-1.7B      & 1.7B & moat, backfires (L26) \\
\bottomrule
\end{tabular}}
\end{table}

\subsection{Corpora and Concepts}

Concept codes are fitted on small in-distribution corpora and evaluated on held-out prompts in three styles: descriptive scenes ($n=100$ per concept), narrative past ($n=100$ per concept), and single-token stubs. The primary concept set is fire, metal, river, and storm, extended with memory and dream for the associative-memory experiments. Corpus robustness is evaluated explicitly in loops L16 to L19 rather than assumed.

\subsection{Seeding and Independence}

Every experiment pre-registers its seeds. Following the L9 stream repair, each generation receives an independent seed derived from a hash of the registered seed, the arm, the concept, and the stub, so arms are compared across genuinely independent sampling streams. Results predating the repair are marked as superseded where later re-runs exist (L18 re-established the dose grid; L19 confirmed the supersession offset at three seeds and two amplitudes). Confirm-tier claims require at least three seeds; the core claim is stated at six.

\subsection{Tiers, Gates, and Statistics}

Experiments run at three tiers: smoke (functional checks, minutes), screen (single- or few-seed effect estimates), and confirm (multi-seed pre-registered criteria). Significance uses permutation tests against shuffled-code or shuffled-label nulls with $10^4$ resamples where distributional assumptions would be doubtful, sign tests across seeds for direction consistency, and registered margins rather than post-hoc thresholds. One early metric was recalibrated after its null was found to be structurally offset: union-top-$K$ rank correlation carries a null floor near $-0.72$, and all reported loadability figures use model-top-$K$ support, whose null is approximately zero, with permutation gates. Where a gate reports an exploratory sweep (the readback threshold search of L14), the sweep size and the absence of multiple-comparison correction are recorded in the gate and restated here (Section~\ref{sec:results:readback}).

\subsection{Reproducibility}

Runs execute in a pinned container (Python 3.10, PyTorch 2.4, aarch64 build with flash attention) with configurations mounted from the versioned \texttt{configs/experiments/} tree (22 experiment configurations). Every claim in Section~\ref{sec:results} cites a gate record in the public ledger by identifier; Appendix~\ref{app:ledger} tabulates the complete ledger, and Appendix~\ref{app:schema} documents the gate schema with a verbatim example. Figures are generated by a script that reads only committed gate records, so each plotted coordinate is traceable to a ledger entry (Appendix~\ref{app:figures}).
